\begin{lemma}[Powers of the basic quadratic phase]\label{mod:finitefield-q0powers}
Let $k$ be a finite field of odd characteristic $p$, let
$\psi:k\to\mathbf C^\times$ be a nontrivial additive character, and
write $\nu:k\to\mathbf C$ for the quadratic character.  For the basic
normalized quadratic phase
\[
 q_0=Q_\psi(1,0)
\]
one has the exact identities
\[
 q_0^2=\nu(-1),\qquad
 q_0^{p-1}=\nu(-1),\qquad
 q_0^p=q_0\nu(-1).
\]
In Lean, $p$ is \texttt{ringChar k} and all three identities are the
three conjuncts of \texttt{quadraticPhase\_powers}.
\LeanFile{LanglandsFirstMainLemma/FiniteField/QuadraticPhasePowers.lean}
\lean{LanglandsFirstMainLemma.quadraticPhase_powers}
\leanok
\SourceText{Lemma lem:q0-powers.}
\uses{mod:finitefield-quadraticphase}
\FormalizationContract
\end{lemma}
\begin{proof}
The quadratic-phase API gives $q_0^2=\nu(-1)$, while the quadratic
character has square one at the nonzero element $-1$.  Hence
$q_0^4=1$.  Since $p$ is an odd prime, either $p\equiv1\pmod4$ or
$p\equiv3\pmod4$.  In the first case, writing $|k|=p^f$ gives
$|k|\equiv1\pmod4$, so the finite-field formula for the quadratic
character at $-1$ gives $\nu(-1)=1$; also $p-1\equiv0\pmod4$, and
therefore $q_0^{p-1}=1=\nu(-1)$.  In the second case,
$p-1\equiv2\pmod4$, whence $q_0^{p-1}=q_0^2=\nu(-1)$.  Finally,
$q_0^p=q_0^{p-1}q_0=q_0\nu(-1)$.
\end{proof}
